\chapter{Libraries of C}\label{ch-c-libraries}

Most of the code a computer runs is C, or is reached through C: the
library of the system, zlib, SQLite, a vendor's driver, and the code a
board's firmware already has. Go and Rust can both make a library
that C can call. Since 0.6 a Zeno program can call any of them, as
they are:

\begin{zeno}[file=maths.zeno, dir=maths]
import c

let m = c.library("m", ["double hypot(double x, double y);", "double frexp(double x, int *exponent);"])
let exponent = [0]
print(m.hypot(3.0, 4.0), m.frexp(8.0, exponent), exponent)
\end{zeno}

\verb|c.library| takes two things: the library, and the functions of
it that the program uses, declared as C declares them. The
declarations are the lines a header has, copied as they are. From
there, each function is called as a function of Zeno is.

This is not the bridge of the last chapter. A value of Python stays
in Python, and Zeno asks for each field of it. C has no such values:
a function of C takes numbers, text and pointers, and gives one back.
So the checker reads the declarations and checks every call before the
program runs, and \verb|zeno build| calls the function directly, with
nothing in between. That is what firmware needs.

A few functions of C can also be written in the \verb|.zeno| file
itself, in an \verb|inline c| block, which gives the same kind of
library with nothing declared. That is the next chapter
(\hyperref[ch-inline]{Another Language in a Zeno File}).

\begin{note}
Nothing changes for a program that does not say \verb|import c|. What
C needs is started by the first \verb|c.library| that runs, and a
program that imports \verb|c| and never calls it starts nothing.
\end{note}

\section{What it needs}

A C compiler: \verb|cc|, \verb|gcc| or \verb|clang|, the first that is
installed, or the one \verb|CC| names, as for \verb|zeno build|. Under
\verb|zeno run|, the first time a program reaches a \verb|c.library|
Zeno writes a small file of C for the functions it declares and builds
it with the library. That takes about half a second, and is kept: the
next run starts at once. It is built again when the declarations
change, or a file of the library does.

The library itself has to be there with what building with it needs.
On Debian and Ubuntu that is the package whose name ends in
\verb|-dev|: \verb|zlib1g-dev| for zlib, \verb|libsqlite3-dev| for
SQLite.

\section{Naming the library}

The first thing given is one text, or a list of them:

\begin{center}
\begin{tabular}{lP{8.6cm}}
\toprule
Given & What it is \\
\midrule
\verb|"z"| & a library of the system, by its name: \verb|-lz|, as C is
  built with it \\
\verb|"crc.c"| & a file of C, built with the rest \\
\verb|"startup.s"| & a file of assembly, which the C compiler assembles
  (\verb|.s|, \verb|.S|) \\
\verb|"libtemp.so"| & a library that was built already: \verb|.so|,
  \verb|.a|, \verb|.o|, \verb|.dylib|, \verb|.dll| \\
\verb|"-DSTEP=3"| & anything that starts with \verb|-| goes to the
  compiler as it is \\
\verb|none| & nothing more: for the C library every program has
  (\verb|strlen|, \verb|abs|) \\
\bottomrule
\end{tabular}
\end{center}

A file is found from the folder of the program, wherever the program
is started from, and so is a header it includes.

\begin{zeno}[file=libc.zeno, dir=maths]
import c

let libc = c.library(none, [
    "size_t strlen(const char *text);",
    "int abs(int n);",
    "char *getenv(const char *name);",
])
print(libc.strlen("zeno"), libc.abs(-7), libc.getenv("NO_SUCH_NAME"))
\end{zeno}

A library that is not there is a \verb|FileError|, which a program can
catch. The message asks if it is installed, with what building with it
needs.

\section{Declaring the functions}

The second thing given is the declarations: one text, or a list of
texts, each declaration ending with \verb|;|. They are C, and only the
functions the program uses are needed.

\begin{itemize}
\item A line that starts with \verb|#|, such as
  \verb|#include <zlib.h>|, is given to C as it is. With the header
  included, the compiler checks each declaration against it, and a
  type the header names (\verb|uLong|, \verb|Bytef|) is known.
\item A \verb|typedef|, a \verb|struct|, a \verb|union| or an
  \verb|enum| is given to C as it is too.
\item \verb|int8_t| to \verb|uint64_t|, \verb|size_t| and \verb|bool|
  can be used with no \verb|#include|.
\item A declaration with no \verb|(| is a value: a constant or a macro
  of a header, a constant of an \verb|enum|, or a variable. It has a
  section below.
\end{itemize}

When the library and its declarations are written out in the program
as texts, which is nearly always, Zeno knows every function before the
program runs, and checks each call as it checks a call of a function
of Zeno:

\begin{zeno}[file=wrong.zeno, dir=maths, cmd=zeno check {file}, fails=yes]
import c

let m = c.library("m", ["double hypot(double x, double y);"])
print(m.hypot("3", 4.0))
print(m.cube(3.0))
\end{zeno}

A type that only a header names (\verb|uLong|,
\texttt{HAL\_Status\-TypeDef}) is taken to be a whole number, as nearly
all such names are, unless it is one of the names headers give floats
(\verb|float32_t|). The C compiler sees the real type, and says so if
the two disagree. A declaration that names a function the library does
not have is an error when \verb|c.library| is called, and not at the
first call.

\begin{why}
C has no way to ask a library what its functions take: a library
knows the names of its functions and nothing of their types. That is
why the program declares them. And Zeno does not read a header by
itself, because a header is written for a C compiler: it is full of
macros, and means something different on each machine. So the
declarations are handed to the C compiler, which is there already
wherever \verb|zeno build| is used, and it works out what
\verb|uLong| or \verb|size_t| is on this machine. (Decisions 44.1 and
44.3.)
\end{why}

\section{Values go across}

\begin{center}
\begin{tabular}{P{5.2cm}P{7.2cm}}
\toprule
C & Zeno \\
\midrule
\verb|int|, \verb|uint8_t|, \verb|long|, \verb|size_t| and their like &
  \verb|int|, checked to fit: an \verb|OverflowError| if it does not \\
\verb|uint64_t| above the top of \verb|int| & \verb|u64| \\
\verb|float|, \verb|double| & \verb|float| \\
\verb|bool|, \verb|_Bool| & \verb|bool| \\
\verb|const char *| given & a \verb|str|, lent for the call, with the
  \verb|\0| C expects \\
\verb|char *| given back & a \verb|str or none|, copied \\
a pointer to numbers, given & a list, lent for the call: copied in,
  and back after \\
any other pointer given back & a C pointer \verb|or none|: given back
  to C, never looked inside \\
\verb|NULL| & \verb|none| \\
\verb|void| & \verb|none| \\
\bottomrule
\end{tabular}
\end{center}

\textbf{A list for a pointer is lent.} Its items are copied into
memory of the type the pointer is to, the function is called, and the
items are copied back into the list afterwards, so a function that
fills a buffer fills the list. That is how \verb|frexp| gave its
exponent in the first program. A pointer to \verb|const| is not copied
back: the function promised not to change it.

\textbf{A pointer that C gives back is kept as it is.} It can be given
back to C, compared with \verb|==|, tested for \verb|none| and printed,
and that is all. This is how a library that hands out a handle
(\verb|FILE *|, \verb|sqlite3 *|, a driver's context) is used. A
function that gives a handle back through a pointer to a pointer
(\verb|sqlite3 **db|) is given a list of one \verb|none|, and finds
the handle in it afterwards.

Zeno never frees memory that C gave it. A library that says the caller
frees what it returns has a function to do it, and the program calls
that, as it would in C.

\section{A library of your own}

Here is a small driver as a board's firmware might have it: a header
and a file of C. Both lie beside the program.

\begin{zeno}[file=sensor.h, dir=sensor, run=no]
#include <stddef.h>
#include <stdint.h>

#define SENSOR_CHANNELS 4

typedef enum { SENSOR_OK, SENSOR_BUSY } sensor_status;
typedef struct sensor sensor;

sensor *sensor_open(uint8_t channel);
sensor_status sensor_read(sensor *s, uint16_t *samples, size_t count);
uint16_t sensor_crc(const uint8_t *data, size_t length);
void sensor_close(sensor *s);
\end{zeno}

\begin{zeno}[file=sensor.c, dir=sensor, run=no]
#include "sensor.h"
#include <stdlib.h>

struct sensor { uint8_t channel; uint16_t next; };

sensor *sensor_open(uint8_t channel) {
    if (channel >= SENSOR_CHANNELS) return NULL;
    sensor *s = malloc(sizeof *s);
    s->channel = channel;
    s->next = 0;
    return s;
}

sensor_status sensor_read(sensor *s, uint16_t *samples, size_t count) {
    for (size_t i = 0; i < count; i++) {
        samples[i] = (uint16_t)(s->channel * 100 + s->next++);
    }
    return SENSOR_OK;
}

uint16_t sensor_crc(const uint8_t *data, size_t length) {
    uint16_t crc = 0xFFFF;
    for (size_t i = 0; i < length; i++) {
        crc ^= data[i];
        for (int bit = 0; bit < 8; bit++) {
            crc = (crc & 1) ? (uint16_t)((crc >> 1) ^ 0xA001) : (uint16_t)(crc >> 1);
        }
    }
    return crc;
}

void sensor_close(sensor *s) { free(s); }
\end{zeno}

\begin{zeno}[file=read.zeno, dir=sensor]
import c

let sensor = c.library("sensor.c", [
    "#include \"sensor.h\"",
    "sensor *sensor_open(uint8_t channel);",
    "sensor_status sensor_read(sensor *s, uint16_t *samples, size_t count);",
    "uint16_t sensor_crc(const uint8_t *data, size_t length);",
    "void sensor_close(sensor *s);",
    "int SENSOR_CHANNELS;",
    "sensor_status SENSOR_OK;",
])

print(sensor.SENSOR_CHANNELS, sensor.sensor_open(9))
let s = sensor.sensor_open(2)
let samples = [0] * 4
if sensor.sensor_read(s, samples, len(samples)) == sensor.SENSOR_OK {
    print(samples)
}
sensor.sensor_close(s)

let frame: list[u8] = [1, 3, 0, 107, 0, 3]
print(sensor.sensor_crc(frame, len(frame)))
\end{zeno}

\verb|sensor_open(9)| gave \verb|NULL|, which is \verb|none|. \verb|s|
is the handle: Zeno never looks inside it, and hands it back.
\verb|samples| was lent and filled. \verb|frame| is a
\verb|list[u8]| (\hyperref[ch-sized-numbers]{Sized Numbers}), which is
what bytes are in Zeno.

\section{Values a header names}

The last two declarations of that program have no \verb|(|, and each
names a value. \verb|SENSOR_CHANNELS| is a macro, and \verb|SENSOR_OK|
a constant of an \verb|enum|. Each is read by its name, as C reads it,
every time the program reads it. C cannot declare a macro, so the
value is not declared again: the header has it, and the declaration
only tells Zeno its type.

That is how the names of a vendor's HAL are used:

\begin{zenofrag}
let hal = c.library("hal.c", [
    "#include \"stm32f4xx_hal.h\"",
    "void HAL_GPIO_WritePin(GPIO_TypeDef *port, uint16_t pin, GPIO_PinState state);",
    "GPIO_TypeDef *GPIOA;",
    "uint16_t GPIO_PIN_5;",
    "GPIO_PinState GPIO_PIN_SET;",
])
hal.HAL_GPIO_WritePin(hal.GPIOA, hal.GPIO_PIN_5, hal.GPIO_PIN_SET)
\end{zenofrag}

A variable that only the library's own file of C has, with no header,
is declared with \verb|extern|, as C would declare it:
\verb|"extern uint32_t uptime_ms;"|.

\section{zeno build: C calls C}

\verb|zeno build| writes C, so a function of a library is called
there as C calls it, with no helper beside the program. The library
is built in with the program: a file of C is compiled with it, and a
library is linked.

\begin{zeno}[file=read.zeno, dir=sensor, show=no, cmd=zeno build {file} ;; ./read ;; zeno build --no-collector --minimal {file} ;; ./read ;; grep -c "" read.c ;; grep "sensor_close(" read.c]
import c

let sensor = c.library("sensor.c", [
    "#include \"sensor.h\"",
    "sensor *sensor_open(uint8_t channel);",
    "sensor_status sensor_read(sensor *s, uint16_t *samples, size_t count);",
    "uint16_t sensor_crc(const uint8_t *data, size_t length);",
    "void sensor_close(sensor *s);",
    "int SENSOR_CHANNELS;",
    "sensor_status SENSOR_OK;",
])

print(sensor.SENSOR_CHANNELS, sensor.sensor_open(9))
let s = sensor.sensor_open(2)
let samples = [0] * 4
if sensor.sensor_read(s, samples, len(samples)) == sensor.SENSOR_OK {
    print(samples)
}
sensor.sensor_close(s)

let frame: list[u8] = [1, 3, 0, 107, 0, 3]
print(sensor.sensor_crc(frame, len(frame)))
\end{zeno}

Every build does it: a plain one, \verb|--firmware|,
\verb|--no-collector|, \verb|--owned| and \verb|--minimal|
(\hyperref[ch-compiling]{Compiling}). The last lines above are the
whole of the minimal build's C counted, and the two lines of it that
have \verb|sensor_close|: the declaration, as the program gave it, and
the call.

A built program prints what \verb|zeno run| prints, and a number that
does not fit stops it with the same \verb|OverflowError|:

\begin{zeno}[file=channel.zeno, dir=sensor, fails=yes, cmd=zeno run {file} ;; zeno build {file} ;; ./channel]
import c

let sensor = c.library("sensor.c", ["#include \"sensor.h\"", "sensor *sensor_open(uint8_t channel);"])
for channel in [2, 300] {
    print(sensor.sensor_open(channel) is none)
}
\end{zeno}

A \verb|--minimal| build does not check this, as it checks nothing.

Three things to know about a build:

\begin{itemize}
\item \textbf{A \texttt{list[u8]} is handed over where it is}, with no
  copy, to a function that takes bytes, and what C writes there is
  what the list holds afterwards. Any other list is copied in and
  back, as under \verb|zeno run|.
\item \textbf{The library and its declarations must be written out}
  in the program as texts, so that they are known while it builds.
\item \textbf{A file of C cannot have the name the build gives the
  program's own C}: \verb|read.c| for \verb|read.zeno|. The build says
  so, and writes nothing.
\end{itemize}

\verb|zeno build --llvm| and \verb|--target| do not take
\verb|import c| yet.

\begin{note}
This was used on Linux x86-64: the system's maths library, zlib,
SQLite, a file of C beside the program, and libraries of Go and Rust.
The tests also run it on Windows (MinGW gcc) and on macOS (Apple
clang), with a file of C and the maths library. It has not been run
with a cross-compiler or on a board.
\end{note}

\section{Libraries of Go and Rust}

Go and Rust each build a library that C can call, and \verb|c.library|
takes that library. The build is done by their own tools.

\textbf{Rust}: mark each function \verb|#[no_mangle] pub extern "C"|,
and build it as a \verb|cdylib| (a \verb|.so|) or a \verb|staticlib|
(a \verb|.a|). With Cargo, that is \verb|crate-type = ["cdylib"]|
under \verb|[lib]| in \verb|Cargo.toml|.

\begin{zeno}[file=filter.rs, dir=rust, run=no]
#[no_mangle]
pub extern "C" fn average(values: *const f64, count: usize) -> f64 {
    let values = unsafe { std::slice::from_raw_parts(values, count) };
    values.iter().sum::<f64>() / count as f64
}
\end{zeno}

\begin{zeno}[file=average.zeno, dir=rust, cmd=rustc -O --crate-type cdylib filter.rs -o libfilter.so ;; zeno run {file}]
import c

let filter = c.library("libfilter.so", "double average(const double *values, size_t count);")
let readings = [20.0, 22.0, 21.0, 25.0]
print(filter.average(readings, len(readings)))
\end{zeno}

\textbf{Go}: mark each function \verb|//export|, and build with the
mode \verb|c-shared| (a \verb|.so|) or \verb|c-archive| (a \verb|.a|).
Go writes \verb|libtemp.h| beside the library, with the declarations.

\begin{zeno}[file=temp.go, dir=go, run=no]
package main

import "C"

//export Celsius
func Celsius(f float64) float64 { return (f - 32) * 5 / 9 }

func main() {}
\end{zeno}

\begin{zeno}[file=temp.zeno, dir=go, cmd=go build -buildmode=c-shared -o libtemp.so temp.go ;; zeno run {file}]
import c

let temp = c.library("libtemp.so", "double Celsius(double f);")
print(temp.Celsius(98.6))
\end{zeno}

There is no \verb|import cpp|. A library of C++ is used through C, as
Go and Rust are: one with \verb|extern "C"| functions is a library of
C (decision 44.12).

\section{When C goes wrong}

A mistake Zeno can see before the program runs is found then: a wrong
argument, a function that is not declared, a declaration Zeno cannot
take, such as a function with \verb|...|. A number that does not fit
is found as the program runs.

A mistake inside C, such as reading memory that is not there, is
another matter. Under \verb|zeno run| C runs beside Zeno, so the
program stops with an error at the line that called it, and Zeno
itself does not crash; \verb|try| cannot catch it. In a built program
C is the program, and it stops as a program of C does.

Text that C gives back is bytes. A byte there that is not part of a
character of UTF-8 is shown as \verb|\u{fffd}| by \verb|zeno run|; a
built program prints the bytes as they are.

\section{What is not there yet}

\begin{itemize}
\item \textbf{The address of a variable} (\verb|&huart2|, a handle of
  a HAL that the vendor's code defines). Write a function of C that
  gives it:\\
  \verb|UART_HandleTypeDef *uart2(void) { return &huart2; }|
\item \textbf{A function of Zeno given to C}, as a callback.
\item \textbf{A struct given or given back whole}, not through a
  pointer.
\item \textbf{A function that takes \texttt{...}}, such as
  \verb|printf|.
\item \textbf{Reading the declarations from a header by itself.} They
  are copied by hand for now, and only the ones used are needed.
\item \textbf{A function of a library kept as a value} is refused by a
  build: \verb|let f = m.hypot| runs under \verb|zeno run| only.
\item \textbf{\texttt{zeno build --llvm}}, and the other languages of
  \verb|--target|.
\end{itemize}

\section{How it is built}

Under \verb|zeno run|, Zeno reads the declarations and writes a small
file of C with a function for each, which takes its arguments from a
table of numbers and calls the real one. The C compiler builds that
with the library. The result is loaded by a small helper written in
C, which runs beside the program and is spoken to as the helpers of
the other languages are (\hyperref[ch-other-languages]{Libraries of
Other Languages}). A fault in C stops the helper and not Zeno.

Under \verb|zeno build| there is no helper. The declarations go into
the program's C as they were written, each call is a call, and a
number is checked to fit before it is given. The example
\verb|examples/modbus.zeno| makes a request of Modbus RTU with the
checksum written in C (\verb|examples/crc.c|), and it is built for
both of the emulated boards with the rest of the examples
(\hyperref[ch-board]{The Board}). The reasons are chapter 44 of
\verb|docs/decisions.md|.

\section{Things to try}

\begin{enumerate}
\item Run the first program of this chapter. Then give \verb|hypot| a
  str and run \verb|zeno check|. Who found the mistake, Zeno or the C
  compiler?
\item Declare \verb|"double cbrt(double x);"| and print
  \verb|m.cbrt(27.0)| under \verb|zeno run|, and then built. The two
  can differ in the last digit. Which one did the C compiler work out
  by itself, before the program ran?
\item Add \verb|uint16_t sensor_peek(const sensor *s);| to the driver
  of this chapter, and call it from Zeno.
\item Give \verb|sensor_read| a list of three floats. What does the
  checker say?
\item Build \verb|read.zeno| with \verb|--minimal| and find the call of
  \verb|sensor_crc| in \verb|read.c|. Where is the list's copy?
\item Call \verb|crc32| of zlib on a file read with
  \verb|files.read_bytes|.
\item Write a function in Go that formats a number with
  \verb|strconv.FormatFloat|, and call it from Zeno. What does Go give
  back, and who frees it?
\end{enumerate}
